\section{Agentic behaviour: SWE-bench Verified and step counts}\label{sec:agentic}

The target workload is a coding agent, so the exploratory period spent four days on SWE-bench Verified
with mini-SWE-agent: 50 instances per build, non-thinking, greedy, 250-step limit. All rows here are
historical\hist.

\begin{table}[t]
\centering\small
\caption{SWE-bench Verified\hist: resolved instances, aggregated from per-instance reports scored on
x86-64. Denominators differ because empty patches and infrastructure errors were excluded case by
case. This is \emph{Verified}, not SWE-bench Pro (on which Qwen reports 61.7 for this model); the two
are not comparable.}\label{tab:swe}
\shrinktofit{\begin{tabular}{@{}lrrl@{}}
\toprule
build & resolved / scored & rate & Wilson 95\,\% \\
\midrule
\IQfour & 38 / 49 & 77.6\,\% & [64.1, 87.0] \\
\Qfive & 38 / 50 & 76.0\,\% & [62.6, 85.7] \\
\Qsix & 37 / 49 & 75.5\,\% & [61.9, 85.4]; bootstrap [63.3, 87.8] \\
\bottomrule
\end{tabular}}
\end{table}

\textbf{The most expensive benchmark ranks the ladder backwards, and that means nothing.}
\Cref{tab:swe} puts the smallest build first. The three builds span 2.1 points; one instance is worth
2 points; one build's interval alone spans 24. A 50-instance agentic benchmark cannot rank neighbouring
quantizations, and an apparent inversion at this size carries no information about them. What the
table does support: all three builds in the 4- to 6-bit range resolve about three quarters of these
instances under this scaffold.

Getting even these numbers right took three generations of totals. Scoring first ran on an ARM64
laptop, and three SWE-bench evaluation images have no ARM64 variant, so specific instances were
silently counted as failures for specific builds. An aggregator then read only the first of seven
concatenated batch reports and published 8 resolved where the true figure was 37. Both were found
by audit and repaired from per-instance records without re-running anything.

\textbf{Below the ladder, a cliff that no cheap instrument saw.} The 3-bit build \Qthree{} was the
fastest configuration of the whole study (116.9\tps{} into an empty cache with MTP depth 8) and scored
84.1 / 81.7 on HumanEval / HumanEval+ with no empty responses: lower than its neighbours, not
alarming. Given agentic tasks, it hit the 250-step limit on \textbf{6 of 6} instances, never converging
once. \Qsix{} on the same six instances converged on all six (mean 45 steps, median 37). I stopped the
3-bit run there, since the remaining 44 instances would have cost about 88 GPU-hours to confirm a
6-to-0 split. This is a qualitative result: $n=6$, one scaffold, one seed, a build that is not part of
the main ladder and whose file is gone. It does not say where between 3 and 4 bits the failure begins.
It is still the sharpest single illustration of the report's theme: every cheap instrument rated this
build acceptable or best, and the one expensive multi-step instrument rated it unusable.

\textbf{Step counts are not a quality proxy in either direction.} Among builds that converge, mean
steps on three shared instances were 23 for \IQfour{} and 29 for \Qfive, against about 45 for the
6-bit build on a different instance set. Per-instance counts overlap and one instance reverses the
order, so nothing is separated at $n=3$. Read with the 3-bit result, the profile across the wider
ladder is non-monotone: the lowest rung never finishes, the middle rungs finish soonest, the top rung
takes longest. Since step count multiplies token cost much more than tokens-per-second does, agentic
efficiency deserves a properly powered study; this was not one.
